\documentclass[tikz,border=8pt]{standalone}
\input{../../../docs/tikz/dag_preamble.tex}

\begin{document}
\begin{tikzpicture}[
  x=1cm,
  y=1cm,
  every node/.append style={outer sep=3pt,font=\sffamily\scriptsize},
  every path/.append style={font=\sffamily\scriptsize}
]

\node[align=left,text width=28cm,font=\sffamily\small,anchor=north west]
  at (0,0) {\textbf{Paper:} Ahorro, flexibilidad de horas e informalidad en el Perú\\
  \textbf{Subtitle:} Un mecanismo laboral en un modelo macroeconómico de agentes heterogéneos\\
  \textbf{Author:} John Svante Barraza Ratachi\\
  \textbf{Source:} merged manuscript, 24 November 2026\\
  \textbf{Scope:} the four stated local assumptions, the informal-hours-share
  definition, and Proposition 1.};

\node[dag_model,text width=5.2cm] (Fixed) at (0,-4.2) {
  \textbf{Assumption 1}\\
  Individual comparison\\
  fixed productivity and returns
};
\node[dag_model,text width=5.2cm] (Preferences) at (7.5,-4.2) {
  \textbf{Assumption 2}\\
  Labor preferences\\
  positive parameters, separable costs
};
\node[dag_model,text width=5.2cm] (Interior) at (15,-4.2) {
  \textbf{Assumption 3}\\
  Regular interior branch\\
  positive, differentiable, slack FOCs
};
\node[dag_model,text width=5.2cm] (Concavity) at (22.5,-4.2) {
  \textbf{Assumption 4}\\
  Local concavity\\
  used only for sign interpretation
};

\node[dag_model,text width=6cm] (Share) at (4.5,-8.6) {
  \textbf{Definition}\\
  Informal-hours share\\
  $s_I=\ell_I/(\ell_F+\ell_I)$
};

\node[dag_result,text width=9cm] (Proposition) at (16.5,-8.6) {
  \textbf{Proposition 1}\\
  Conditional composition of hours\\
  derivative identity and the negative, zero,
  and positive elasticity-order cases
};

\draw[dag_arrow] (Fixed.south) -- ++(0,-1.0) -| ([xshift=-3.3cm]Proposition.north);
\draw[dag_arrow] (Preferences.south) -- ++(0,-1.6) -| ([xshift=-1.2cm]Proposition.north);
\draw[dag_arrow] (Interior.south) -- ([xshift=0.8cm]Proposition.north);
\draw[dag_arrow] (Concavity.south) -- ++(0,-1.0) -| ([xshift=3.3cm]Proposition.north);
\draw[dag_arrow] (Share.east) -- (Proposition.west);

\node[align=left,text width=28cm,font=\sffamily\scriptsize,anchor=north west]
  at (0,-11.8) {\textit{Reading guide:} arrows run from each source condition
  or definition to the result that uses it. Assumptions 1--3 establish the
  identity; Assumption 4 is used only inside the proposition's conditional
  sign clause. The result is local and conditional: no node asserts global HJB
  regularity or concavity, boundary or active-cap behavior, aggregation, or
  general-equilibrium existence.};

\end{tikzpicture}
\end{document}
